% --- 衬线体风格标题页 ---
\pdfbookmark{标题页}{title}
\thispagestyle{empty}

\begin{center}
    \vspace*{\stretch{2.5}}

    % 1. 中文主标题：使用衬线体 (\rmfamily)，加粗
    % 在 CTeX 中，\bfseries 作用于宋体通常会产生“伪粗体”或调用“堆宋”
    {\fontsize{34pt}{40pt}\selectfont\rmfamily\bfseries\color{ThemeWine}
        数学分析讲义（第二册）\\[0.4em] 习题解析
    }

    \vspace{1.2cm}

    % 2. 经典双线装饰 (Wine + Rose)
    % 衬线体配合稍微拉长的装饰线非常有感觉
    {\color{ThemeWine}\hrule height 2.2pt}
    \vspace{2.5pt}
    {\color{ThemeRose}\hrule height 0.8pt}

    \vspace{1.2cm}

    % 3. 英文标题：使用衬线体
    % \textsf 被替换为 \rmfamily，并使用 \itshape (斜体) 增加数学感
    {\fontsize{18pt}{22pt}\selectfont\rmfamily\bfseries\color{ThemeDarkBlue}
        Mathematical Analysis (Volume II)
    } \\
    \vspace{0.6em}
    {\large\rmfamily\bfseries\color{ThemeDarkBlue} Exercise Solutions}

    \vspace*{\stretch{2}}

    % 4. 作者与机构：楷体 (\kaishu) 是中文排版中极佳的衬线补充
    {\Large\kaishu\color{ThemeDarkBlue} 李清臣} \footnote{%
        \color{ThemeDarkBlue}\rmfamily
        闻道有先后，疏漏之处在所难免。如有任何问题，恳请读者在本文档的 GitHub 仓库
        \href{https://github.com/wanzhao-ysy/math-analysis-vol2-solutions}{\texttt{math-analysis-vol2-solutions}}
        提交 Issue，或发送邮件至 \href{mailto:wanzhao_ysy@qq.com}{\texttt{wanzhao\_ysy@qq.com}}。
    }

    \vspace{0.5cm}
    {\large\kaishu\color{ThemeMidBlue} 中国科学技术大学}

    \vspace{2cm}

    % 5. 底部装饰：一个精致的衬线符号
    {\color{ThemeRose}\Large $\mathcal{Q.E.D.}$}

    \vspace{0.8cm}
    {\large\rmfamily\color{ThemeMidBlue}\today}

    \vspace*{\stretch{3.5}}
\end{center}

\newpage
\thispagestyle{empty}

\endinput
